\documentclass[11pt]{article}
\usepackage[a4paper,margin=2.2cm]{geometry}
\usepackage{booktabs}
\usepackage{array}
\usepackage{microtype}
\usepackage[hidelinks]{hyperref}
\title{\vspace{-1.2em}\textbf{ASTRA $\leftarrow$ AlphaEvolve: Performance Impact Assessment}\\[0.2em]
\large Has the post-AlphaEvolve work improved ASTRA, and was it worthwhile?}
\author{ASTRA Engineering Note}
\date{11 July 2026}
\begin{document}
\maketitle
\vspace{-1.5em}

\section*{1. The change, in one paragraph}
After the AlphaEvolve paper was shared as a reference, ASTRA's discovery
pipeline was re-architected from a \textbf{single-shot LLM text generator judged by
LLM-as-judge and regex} (circular, unfalsifiable, and the structural source of
the project's long-standing ``fiction problem'') into an
\textbf{AlphaEvolve-style evolutionary loop with a machine-checkable
\texttt{EVALUATE} on real archival data}. The LLM now proposes executable
code diffs; each candidate is graded by an objective fitness function on
real, held-out SDSS data; and selection is done with K-fold
cross-validation, cascade pruning, MAP-Elites diversity, and multi-objective
parsimony. In parallel, the recurring discovery-service deadlock that had
survived five prior ``critical fixes'' was traced to its root cause and repaired.

\section*{2. Measured impact (verified on real data, not simulated)}
All numbers below are from runs against \emph{real} SDSS photometry, persisted
under \texttt{\textasciitilde/.astra\_persistent/evolved\_programs/} and reproducible.

\begin{center}
\small
\begin{tabular}{l c c c}
\toprule
\textbf{Metric / capability} & \textbf{Before} & \textbf{After} & \textbf{Change} \\
\midrule
Photometric redshift $\sigma_{\mathrm{NMAD}}$ (held-out) & 0.0491 & 0.0199 & $-59\%$ \\
Photo-$z$ outlier rate $\eta$ & 0.036 & 0.005 & $-86\%$ \\
STAR/GAL/QSO balanced accuracy (2nd task) & 0.333 & 0.912 & $+174\%$ \\
Selection$\rightarrow$TEST gap (5-fold CV) & 0.0021 & 0.0014 & $-33\%$ \\
corr(selection, held-out TEST) & $+0.758$ & $+0.815$ & sharper \\
Wall-clock throughput (4 workers) & $1.0\times$ & $1.86\times$ & $+86\%$ \\
\midrule
Discovery judgement & LLM-judge + regex & machine fitness on real data & structural \\
Service init deadlock & recurring hang & \emph{root-cause} fix (DEVNULL) & resolved \\
\bottomrule
\end{tabular}
\end{center}

\textbf{Three results deserve emphasis.} First, on the primary task the evolved
pipelines roughly \textbf{halve} the redshift error ($\sigma_{\mathrm{NMAD}}$
0.0491$\rightarrow$0.0199), with the best program being LLM-authored (raw
magnitudes + four colours $\rightarrow$ regularised GradientBoosting). Second,
the \textbf{same engine generalises} to a different problem type and metric
(three-class classification, 0.333$\rightarrow$0.912 balanced accuracy) --- so
the gain is not photo-$z$-specific. Third, K-fold cross-validation made the
measurement itself \emph{trustworthy}: it shrank the selection$\rightarrow$test
gap by a third and exposed an earlier single-evaluation ``trend'' as noise,
which is exactly the discipline that was missing before.

\section*{3. What ``improved'' actually means here}
The most important change is \emph{not} in the table. Astrophysics has objective
metrics everywhere ($\chi^{2}$, $\sigma_{\mathrm{NMAD}}$, false-alarm
probability, S/N), but ASTRA previously never closed the loop on them ---
discovery quality was asserted, not measured. AlphaEvolve's central lesson is
that the \texttt{EVALUATE} function \emph{is} the system. Giving ASTRA one
converts its output from \textbf{unfalsifiable claims into reproducible,
machine-graded results}. This is the structural fix for the fiction problem the
project's own rules exist to guard against: a candidate that does not reduce
error on real held-out data now simply cannot be promoted as a discovery.

A secondary but concrete win: the watchdog that supervises the always-on
discovery service was spawning its child with \texttt{stdout/stderr=PIPE} and
never draining the pipes, so the child's next log flush blocked forever --- this
was the true cause of the init hang that v5.0--v7.0 each claimed to fix.
Switching to \texttt{DEVNULL} removes it at the root instead of adding another
bandage.

\section*{4. Honest gaps (what is not yet done)}
\begin{itemize}\setlength\itemsep{0.1em}
\item \textbf{Live integration is partial.} The consumer that ingests
machine-verified discoveries into the auto-start loop is wired and
isolated-tested, but the always-on discovery service is currently \emph{stopped};
the offline evolutionary loop is verified, the live path is not yet exercised.
\item \textbf{Code generation bypasses \texttt{system.answer()}} because STAN's
\texttt{answer()} returns canned content that ignores the prompt; the LLM
diff-proposer therefore uses the direct Claude gateway instead. This works but is
not the intended integration point.
\item \textbf{Only one machine-verified discovery has been emitted to date}
(\texttt{evolved\_discoveries.json}). The loop produces them readily; live
ingestion depends on the service running.
\item \textbf{Remaining hardening} is an OS-level sandbox for untrusted generated
code. A third task (period finding) was deferred by decision --- generality was
already proven.
\end{itemize}

\section*{5. Was the work worthwhile?}
\textbf{Yes --- on the merits, clearly.} The work gave ASTRA the single
component AlphaEvolve shows is decisive: an automated, machine-checkable fitness
signal on real data, with honest measurement (error bars, cross-validation)
built in. The headline gains are real, reproducible, and generalisable across two
task types, and a deadlock that had defeated five prior fixes was finally
diagnosed and removed at its root. That moves ASTRA from ``asserts it discovered
something'' to ``can prove it on held-out data'' --- a qualitative change in
trustworthiness, which is precisely what this project's anti-fiction rules have
been demanding.

\textbf{With one honest caveat on timing.} The payoff is currently realised in
the \emph{standalone} evolutionary loop, not yet in the always-on discovery
service. What remains is integration hardening (run the service, route through
the intended entry point, sandbox generated code) --- not redoing the core. So:
the work was worthwhile and is structurally sound; the last mile, not the
foundation, is what is left to complete.

\end{document}
